% ---- score2 standard preamble (from: python3 score2_tools.py preamble) ----
\documentclass[10pt]{extarticle}
\usepackage[a4paper,margin=1.5cm]{geometry}
\usepackage[T1]{fontenc}
\usepackage{amsmath,amssymb}
\usepackage{graphicx}
\usepackage[export]{adjustbox}
\usepackage{multicol}
\usepackage{enumitem}
\usepackage{booktabs,array,tabularx}
\setlength{\columnsep}{0.7cm}
\setlength{\parindent}{0pt}
\setlength{\parskip}{3pt}
\pagestyle{empty}
\makeatletter
% figure registry (figdefs.tex is generated by the extract step)
\newcommand{\qfigdef}[4]{\@namedef{qf@file@#1}{#2}\@namedef{qf@w@#1}{#3}\@namedef{qf@wide@#1}{#4}}
\newcommand{\qf@check}[1]{\@ifundefined{qf@file@#1}{\PackageError{score2}{Unknown figure id #1}{}}{}}
\newcommand{\qf@img}[4][]{\includegraphics[width=#2\dimexpr\@nameuse{qf@w@#4}pt\relax,max width=#3,max height=0.45\textheight,#1]{\@nameuse{qf@file@#4}}}
% Full-width questions: finalize wraps every question that contains a wide figure in \qfullbegin...\qfullend,
% so the WHOLE question (stem, figures, options) is set across the page in reading order, then the two columns
% resume. The question's head - its number, stem and first figure (\qfullsplit marks the end of the head) -
% is kept on one page, so a number is never left at a page foot with its figure overleaf; the rest may break
% across pages normally, so no big gaps are left. Without a \qfullsplit, the whole question is kept together
% if it is under 45% of a page. Never write these yourself.
\newif\ifqf@full
\newif\ifqf@boxed
\newsavebox{\qf@box}
\newcommand{\qf@open}{\global\qf@fulltrue\par\global\qf@boxedtrue
  \setbox\qf@box\vbox\bgroup\hsize\textwidth\linewidth\textwidth}
\newcommand{\qf@emit}[1]{\par\egroup\global\qf@boxedfalse
  \ifdim\dimexpr\ht\qf@box+\dp\qf@box\relax<#1\textheight\noindent\usebox{\qf@box}\par\else\unvbox\qf@box\fi}
\newcommand{\qfullsplit}{\ifqf@boxed\qf@emit{0.85}\fi}
\newcommand{\qf@close}{\ifqf@boxed\qf@emit{0.45}\fi\global\qf@fullfalse}
\newcommand{\qfullbegin}{\end{multicols}\qf@open}
\newcommand{\qfullend}{\qf@close\begin{multicols}{2}}
\newcommand{\qfullstart}{\qf@open}
\newcommand{\qfulldone}{\qf@close}
% \qfig[scale]{id}: block figure, centred, directly where it is written (size set by the tool).
% Fallback if a wide figure is not inside a full-width question: break out of the columns in place - no floats.
\newcommand{\qfig}[2][1]{\qf@check{#2}%
  \ifqf@full\qf@blk{#1}{#2}\else
  \ifnum\@nameuse{qf@wide@#2}=1 \expandafter\@firstoftwo\else\expandafter\@secondoftwo\fi
  {\end{multicols}\begin{center}\qf@img{#1}{\textwidth}{#2}\end{center}\begin{multicols}{2}}%
  {\qf@blk{#1}{#2}}\fi}
\newcommand{\qf@blk}[2]{\par\smallskip{\centering\qf@img{#1}{\linewidth}{#2}\par}\smallskip}
% \qfigi[scale]{id}: inline figure (inside option lists / tables / minipages); never breaks columns.
\newcommand{\qfigi}[2][1]{\qf@check{#2}\qf@img[valign=c]{#1}{\linewidth}{#2}}
\makeatother
% structure macros - lint reads these to check completeness and figure placement
\newcommand{\papersec}[1]{\par\bigskip{\centering\Large\bfseries PAPER-#1\par}\medskip}
\newcommand{\partsec}[2]{\par\medskip\noindent{\large\bfseries PART-#1: #2}\par\smallskip}
\newcommand{\sect}[1]{\par\smallskip\noindent\textbf{#1}\par}
\newcommand{\qn}[1]{\par\smallskip\noindent\textbf{#1.}~}
\newcommand{\paraq}[2]{\par\smallskip\noindent\textbf{Paragraph for Questions #1 and #2}\par}
\newenvironment{qoptions}{\begin{itemize}[leftmargin=1.8em,itemsep=1pt,topsep=2pt,label={}]}{\end{itemize}}
\newcommand{\opt}[2]{\item[(#1)] #2}
% figure registry (generated by score2_tools.py finalize)
\qfigdef{p1_phy_s1_q03_1}{TOAS - Test 01 (30 Jan 2022)-img/p1_phy_s1_q03_1.png}{88.4}{0}
\qfigdef{p1_phy_s2_q06_1}{TOAS - Test 01 (30 Jan 2022)-img/p1_phy_s2_q06_1.png}{164.9}{0}
\qfigdef{p1_phy_s2_q08_1}{TOAS - Test 01 (30 Jan 2022)-img/p1_phy_s2_q08_1.png}{114.6}{0}
\qfigdef{p1_phy_s2_q09_1}{TOAS - Test 01 (30 Jan 2022)-img/p1_phy_s2_q09_1.png}{72.0}{0}
\qfigdef{p1_phy_s2_q10_1}{TOAS - Test 01 (30 Jan 2022)-img/p1_phy_s2_q10_1.png}{117.1}{0}
\qfigdef{p1_phy_s3_para1_2_1}{TOAS - Test 01 (30 Jan 2022)-img/p1_phy_s3_para1_2_1.png}{117.1}{0}
\qfigdef{p1_phy_s3_para5_6_1}{TOAS - Test 01 (30 Jan 2022)-img/p1_phy_s3_para5_6_1.png}{114.8}{0}
\qfigdef{p1_phy_s4_q09_1}{TOAS - Test 01 (30 Jan 2022)-img/p1_phy_s4_q09_1.png}{120.0}{0}
\qfigdef{p1_chem_s1_q03_A}{TOAS - Test 01 (30 Jan 2022)-img/p1_chem_s1_q03_A.png}{89.6}{0}
\qfigdef{p1_chem_s1_q03_B}{TOAS - Test 01 (30 Jan 2022)-img/p1_chem_s1_q03_B.png}{102.3}{0}
\qfigdef{p1_chem_s1_q03_C}{TOAS - Test 01 (30 Jan 2022)-img/p1_chem_s1_q03_C.png}{44.1}{0}
\qfigdef{p1_chem_s1_q03_D}{TOAS - Test 01 (30 Jan 2022)-img/p1_chem_s1_q03_D.png}{51.7}{0}
\qfigdef{p1_chem_s1_q04_A}{TOAS - Test 01 (30 Jan 2022)-img/p1_chem_s1_q04_A.png}{161.5}{0}
\qfigdef{p1_chem_s1_q04_B}{TOAS - Test 01 (30 Jan 2022)-img/p1_chem_s1_q04_B.png}{149.4}{0}
\qfigdef{p1_chem_s1_q04_C}{TOAS - Test 01 (30 Jan 2022)-img/p1_chem_s1_q04_C.png}{147.2}{0}
\qfigdef{p1_chem_s1_q04_D}{TOAS - Test 01 (30 Jan 2022)-img/p1_chem_s1_q04_D.png}{131.8}{0}
\qfigdef{p1_chem_s2_q08_D}{TOAS - Test 01 (30 Jan 2022)-img/p1_chem_s2_q08_D.png}{30.8}{0}
\qfigdef{p1_chem_s2_q09_A}{TOAS - Test 01 (30 Jan 2022)-img/p1_chem_s2_q09_A.png}{98.0}{0}
\qfigdef{p1_chem_s2_q09_B}{TOAS - Test 01 (30 Jan 2022)-img/p1_chem_s2_q09_B.png}{76.8}{0}
\qfigdef{p1_chem_s2_q09_C}{TOAS - Test 01 (30 Jan 2022)-img/p1_chem_s2_q09_C.png}{111.2}{0}
\qfigdef{p1_chem_s2_q09_D}{TOAS - Test 01 (30 Jan 2022)-img/p1_chem_s2_q09_D.png}{76.5}{0}
\qfigdef{p1_chem_s2_q10_1}{TOAS - Test 01 (30 Jan 2022)-img/p1_chem_s2_q10_1.png}{306.6}{1}
\qfigdef{p1_chem_s3_para5_6_1}{TOAS - Test 01 (30 Jan 2022)-img/p1_chem_s3_para5_6_1.png}{150.6}{0}
\qfigdef{p2_phy_s1_q01_1}{TOAS - Test 01 (30 Jan 2022)-img/p2_phy_s1_q01_1.png}{140.1}{0}
\qfigdef{p2_phy_s1_q02_1}{TOAS - Test 01 (30 Jan 2022)-img/p2_phy_s1_q02_1.png}{103.6}{0}
\qfigdef{p2_phy_s1_q04_1}{TOAS - Test 01 (30 Jan 2022)-img/p2_phy_s1_q04_1.png}{151.6}{0}
\qfigdef{p2_phy_s1_q07_1}{TOAS - Test 01 (30 Jan 2022)-img/p2_phy_s1_q07_1.png}{148.2}{0}
\qfigdef{p2_phy_s1_q08_A}{TOAS - Test 01 (30 Jan 2022)-img/p2_phy_s1_q08_A.png}{94.2}{0}
\qfigdef{p2_phy_s1_q10_1}{TOAS - Test 01 (30 Jan 2022)-img/p2_phy_s1_q10_1.png}{105.1}{0}
\qfigdef{p2_phy_s1_q12_A}{TOAS - Test 01 (30 Jan 2022)-img/p2_phy_s1_q12_A.png}{122.4}{0}
\qfigdef{p2_phy_s1_q14_1}{TOAS - Test 01 (30 Jan 2022)-img/p2_phy_s1_q14_1.png}{146.3}{0}
\qfigdef{p2_phy_s1_q15_1}{TOAS - Test 01 (30 Jan 2022)-img/p2_phy_s1_q15_1.png}{97.3}{0}
\qfigdef{p2_phy_s1_q17_1}{TOAS - Test 01 (30 Jan 2022)-img/p2_phy_s1_q17_1.png}{22.0}{0}
\qfigdef{p2_phy_s1_q19_1}{TOAS - Test 01 (30 Jan 2022)-img/p2_phy_s1_q19_1.png}{52.7}{0}
\qfigdef{p2_phy_s1_q20_1}{TOAS - Test 01 (30 Jan 2022)-img/p2_phy_s1_q20_1.png}{84.6}{0}
\qfigdef{p2_phy_s2_q01_1}{TOAS - Test 01 (30 Jan 2022)-img/p2_phy_s2_q01_1.png}{129.6}{0}
\qfigdef{p2_phy_s2_q02_1}{TOAS - Test 01 (30 Jan 2022)-img/p2_phy_s2_q02_1.png}{112.4}{0}
\qfigdef{p2_phy_s2_q04_1}{TOAS - Test 01 (30 Jan 2022)-img/p2_phy_s2_q04_1.png}{34.6}{0}
\qfigdef{p2_phy_s2_q05_1}{TOAS - Test 01 (30 Jan 2022)-img/p2_phy_s2_q05_1.png}{103.6}{0}
\qfigdef{p2_phy_s2_q07_1}{TOAS - Test 01 (30 Jan 2022)-img/p2_phy_s2_q07_1.png}{65.8}{0}
\qfigdef{p2_phy_s2_q09_1}{TOAS - Test 01 (30 Jan 2022)-img/p2_phy_s2_q09_1.png}{90.3}{0}
\qfigdef{p2_phy_s2_q10_1}{TOAS - Test 01 (30 Jan 2022)-img/p2_phy_s2_q10_1.png}{108.4}{0}
\qfigdef{p2_chem_s1_q02_1}{TOAS - Test 01 (30 Jan 2022)-img/p2_chem_s1_q02_1.png}{49.7}{0}
\qfigdef{p2_chem_s1_q04_1}{TOAS - Test 01 (30 Jan 2022)-img/p2_chem_s1_q04_1.png}{54.6}{0}
\qfigdef{p2_chem_s1_q04_A}{TOAS - Test 01 (30 Jan 2022)-img/p2_chem_s1_q04_A.png}{39.7}{0}
\qfigdef{p2_chem_s1_q04_B}{TOAS - Test 01 (30 Jan 2022)-img/p2_chem_s1_q04_B.png}{42.6}{0}
\qfigdef{p2_chem_s1_q04_C}{TOAS - Test 01 (30 Jan 2022)-img/p2_chem_s1_q04_C.png}{39.7}{0}
\qfigdef{p2_chem_s1_q04_D}{TOAS - Test 01 (30 Jan 2022)-img/p2_chem_s1_q04_D.png}{39.9}{0}
\qfigdef{p2_chem_s1_q14_A}{TOAS - Test 01 (30 Jan 2022)-img/p2_chem_s1_q14_A.png}{114.2}{0}
\qfigdef{p2_chem_s1_q14_B}{TOAS - Test 01 (30 Jan 2022)-img/p2_chem_s1_q14_B.png}{105.1}{0}
\qfigdef{p2_chem_s1_q14_C}{TOAS - Test 01 (30 Jan 2022)-img/p2_chem_s1_q14_C.png}{114.6}{0}
\qfigdef{p2_chem_s1_q14_D}{TOAS - Test 01 (30 Jan 2022)-img/p2_chem_s1_q14_D.png}{133.9}{0}
\qfigdef{p2_chem_s1_q15_1}{TOAS - Test 01 (30 Jan 2022)-img/p2_chem_s1_q15_1.png}{106.1}{0}
\qfigdef{p2_chem_s1_q16_B}{TOAS - Test 01 (30 Jan 2022)-img/p2_chem_s1_q16_B.png}{58.0}{0}
\qfigdef{p2_chem_s1_q16_2}{TOAS - Test 01 (30 Jan 2022)-img/p2_chem_s1_q16_2.png}{43.0}{0}
\qfigdef{p2_chem_s1_q17_A}{TOAS - Test 01 (30 Jan 2022)-img/p2_chem_s1_q17_A.png}{58.8}{0}
\qfigdef{p2_chem_s1_q17_B}{TOAS - Test 01 (30 Jan 2022)-img/p2_chem_s1_q17_B.png}{27.4}{0}
\qfigdef{p2_chem_s1_q17_C}{TOAS - Test 01 (30 Jan 2022)-img/p2_chem_s1_q17_C.png}{49.2}{0}
\qfigdef{p2_chem_s1_q17_D}{TOAS - Test 01 (30 Jan 2022)-img/p2_chem_s1_q17_D.png}{55.1}{0}
\qfigdef{p2_chem_s1_q18_1}{TOAS - Test 01 (30 Jan 2022)-img/p2_chem_s1_q18_1.png}{168.2}{0}
\qfigdef{p2_chem_s1_q19_B}{TOAS - Test 01 (30 Jan 2022)-img/p2_chem_s1_q19_B.png}{137.3}{0}
\qfigdef{p2_chem_s1_q20_1}{TOAS - Test 01 (30 Jan 2022)-img/p2_chem_s1_q20_1.png}{51.2}{0}
\qfigdef{p2_chem_s1_q20_A}{TOAS - Test 01 (30 Jan 2022)-img/p2_chem_s1_q20_A.png}{52.9}{0}
\qfigdef{p2_chem_s1_q20_B}{TOAS - Test 01 (30 Jan 2022)-img/p2_chem_s1_q20_B.png}{50.7}{0}
\qfigdef{p2_chem_s1_q20_C}{TOAS - Test 01 (30 Jan 2022)-img/p2_chem_s1_q20_C.png}{51.9}{0}
\qfigdef{p2_chem_s1_q20_D}{TOAS - Test 01 (30 Jan 2022)-img/p2_chem_s1_q20_D.png}{51.1}{0}
\qfigdef{p2_chem_s2_q08_1}{TOAS - Test 01 (30 Jan 2022)-img/p2_chem_s2_q08_1.png}{144.9}{0}
\qfigdef{p2_chem_s2_q08_2}{TOAS - Test 01 (30 Jan 2022)-img/p2_chem_s2_q08_2.png}{62.0}{0}
\qfigdef{p2_chem_s2_q09_1}{TOAS - Test 01 (30 Jan 2022)-img/p2_chem_s2_q09_1.png}{150.1}{0}
\qfigdef{p2_chem_s2_q10_1}{TOAS - Test 01 (30 Jan 2022)-img/p2_chem_s2_q10_1.png}{150.1}{0}
\qfigdef{p2_math_s1_q01_1}{TOAS - Test 01 (30 Jan 2022)-img/p2_math_s1_q01_1.png}{79.6}{0}
\qfigdef{p2_math_s1_q05_1}{TOAS - Test 01 (30 Jan 2022)-img/p2_math_s1_q05_1.png}{77.9}{0}
\qfigdef{p2_math_s1_q20_A}{TOAS - Test 01 (30 Jan 2022)-img/p2_math_s1_q20_A.png}{92.0}{0}
\qfigdef{p2_math_s1_q20_B}{TOAS - Test 01 (30 Jan 2022)-img/p2_math_s1_q20_B.png}{92.8}{0}
\qfigdef{p2_math_s1_q20_C}{TOAS - Test 01 (30 Jan 2022)-img/p2_math_s1_q20_C.png}{96.0}{0}
\qfigdef{p2_math_s1_q20_D}{TOAS - Test 01 (30 Jan 2022)-img/p2_math_s1_q20_D.png}{93.0}{0}
\raggedcolumns
\makeatletter\@beginparpenalty=10000\makeatother
\newcolumntype{L}[1]{>{\raggedright\arraybackslash}p{#1}}
\begin{document}
\begin{center}{\Large\bfseries ALLEN TOAS (Phase-TOAS, TNAS \& TAAS) --- Test 01 (30 Jan 2022)}\end{center}
\begin{multicols}{2}
\papersec{1}
\partsec{1}{PHYSICS}
\sect{SECTION-I(i)}
\qn{1} The maximum force on a body of mass $m$ executing SHM is $p$ and the maximum kinetic energy of the body is $q$. What is its time period?
\begin{qoptions}
\opt{A}{$\dfrac{2\pi}{p}\sqrt{2mq}$}
\opt{B}{$\dfrac{2\pi}{q}\sqrt{2mp}$}
\opt{C}{$\dfrac{\pi}{p}\sqrt{2mq}$}
\opt{D}{$\dfrac{\pi}{q}\sqrt{2mp}$}
\end{qoptions}
\qn{2} Potential energy of a particle executing SHM varies sinusoidally. If the frequency of oscillation of the particle is $f$ then frequency of oscillation of potential energy is
\begin{qoptions}
\opt{A}{$f$}
\opt{B}{$\dfrac{f}{2}$}
\opt{C}{$4f$}
\opt{D}{$2f$}
\end{qoptions}
\qn{3} A liquid of density $\rho$ is filled in a conical vessel as shown in figure. Force exerted by liquid on side wall is :
\qfig{p1_phy_s1_q03_1}
\begin{qoptions}
\opt{A}{$\dfrac{2}{3}\pi R^2\rho gh$}
\opt{B}{$\dfrac{1}{3}\pi R^2\rho gh$}
\opt{C}{$\pi R^2\rho gh$}
\opt{D}{$0$}
\end{qoptions}
\qn{4} A small ball of density $4\rho_0$ is released from rest just below the surface of a liquid. The density of liquid increases with depth as $\rho=\rho_0(1+ay)$ where $a=2\,\mathrm{m^{-1}}$ is a constant. Find the time period of its oscillation. (neglect the viscosity effects)
\begin{qoptions}
\opt{A}{$\dfrac{2\pi}{\sqrt5}$ sec}
\opt{B}{$\dfrac{\pi}{\sqrt5}$ sec}
\opt{C}{$\dfrac{\pi}{2\sqrt5}$ sec}
\opt{D}{$\dfrac{3\pi}{2\sqrt5}$ sec}
\end{qoptions}
\sect{SECTION-I(ii)}
\qn{5} In a simple harmonic motion of amplitude $A$ (symbols have their usual meaning)
\begin{qoptions}
\opt{A}{Time taken by particle to move from mean position to $\dfrac{A}{2}$ is equal to time taken by particle to move from $\dfrac{A}{2}$ to $A$}
\opt{B}{Average velocity of particle in one time period $T$ is $\dfrac{4A}{T}$}
\opt{C}{Net work done on block in one time period $T$ is $0$}
\opt{D}{Particle has maximum kinetic energy at mean position}
\end{qoptions}
\qn{6} Two spring, each of outstretched length 20 cm but having different spring constants $k_1=1000$ N/m and $k_2=3000$ N/m are attached to two opposite faces of a block of mass $m=100$ g kept on a smooth horizontal surface as shown in figure.
\qfig{p1_phy_s2_q06_1}
The outer ends of the two springs are now attached to two pins $P_1$ and $P_2$ whose locations are shown in the figure, as a result of this, the block acquired a new equilibrium position. The block has been displaced by small amount from its equilibrium position and released to perform simple harmonic motion, then.
\begin{qoptions}
\opt{A}{New equilibrium position is at 35 cm from $P_1$ and time period of simple harmonic motion is $\dfrac{\pi}{100}$ s.}
\opt{B}{New equilibrium position is at 20 cm from $P_1$ and time period of simple harmonic motion is $\dfrac{\pi}{100}$ s.}
\opt{C}{New equilibrium position is at 35 cm from $P_1$ and time period of simple harmonic motion is $\dfrac{\pi}{26}$ s.}
\opt{D}{New equilibrium position is at 30 cm from $P_1$ and time period of simple harmonic is $\dfrac{\pi}{26}$ s}
\end{qoptions}
\qn{7} A particle of mass $m$ moves along X-axis. If $v$ is the velocity at a distance $x$ from a origin on the X-axis such that $v^2=a-bx^2$, where $a$ and $b$ are constant, then (Assume that potential energy of the particle is zero when it is at origin)
\begin{qoptions}
\opt{A}{the particle continues to move along the positive x-direction only}
\opt{B}{the motion is simple harmonic}
\opt{C}{the particle oscillates with a frequency equal to $\dfrac{\sqrt b}{2\pi}$}
\opt{D}{the total energy of the particle is $\dfrac{ma}{2}$}
\end{qoptions}
\qn{8} Water is flowing on a horizontal surface where velocity varies with height $y$ from the ground as $v=\dfrac{v_0}{H^2}\left(2y^2-\dfrac{y^3}{H}\right)$ where $H$ is the depth of water. Coefficient of viscosity of water is $\eta$. Then choose the correct option(s).
\qfig{p1_phy_s2_q08_1}
\begin{qoptions}
\opt{A}{Shear stress at $y=\dfrac{H}{2}$ is $\dfrac{5v_0\eta}{4H}$}
\opt{B}{Shear stress at $y=H$ is $\dfrac{\eta v_0}{H}$}
\opt{C}{Viscous force acting on the boat of base area $A$ is $\dfrac{\eta A v_0}{H}$}
\opt{D}{Viscous force acting on the boat of base area $A$ is $\dfrac{\eta A v_0}{2H}$}
\end{qoptions}
\qn{9} A sphere of radius $r$ and density $\rho$ is released from the surface of a liquid of density $\dfrac{\rho}{2}$ and coefficient of viscosity $\eta$. It achieves terminal velocity just before reaching the ground. Then choose the correct option(s).
\qfig{p1_phy_s2_q09_1}
\begin{qoptions}
\opt{A}{Terminal velocity is $\dfrac{2r^2\rho g}{9\eta}$}
\opt{B}{Terminal velocity is $\dfrac{r^2\rho g}{9\eta}$}
\opt{C}{Magnitude of work done by viscous force on the sphere is $\dfrac{2\pi r^3\rho g}{3}\left(H-\dfrac{\rho^2r^4g}{81\eta^2}\right)$}
\opt{D}{Magnitude of work done by viscous force on the sphere is $\dfrac43\pi r^3\rho g\left(H-\dfrac{\rho^2r^4g}{9\eta^2}\right)$}
\end{qoptions}
\qn{10} A small solid ball of density $r$ is held inside at point $A$ in a closed cubical container of side $L$, filled with an ideal liquid of density $4\rho$ as shown in the figure. Now, if the container starts moving with constant acceleration $a$ horizontally and the ball is released from point $A$ simultaneously, then
\qfig{p1_phy_s2_q10_1}
\begin{qoptions}
\opt{A}{For ball to hit the top of container at end $Q$, $a=3g$}
\opt{B}{For ball to hit the top of container at end $Q$, $a=2g$}
\opt{C}{Ball hits the top of container at end $Q$ after a time $t=\sqrt{\dfrac{L}{3g}}$}
\opt{D}{Ball hits the top of container at end $Q$ after a time $t=\sqrt{\dfrac{2L}{3g}}$}
\end{qoptions}
\sect{SECTION-II(i)}
\paraq{1}{2}
A hemispherical tank of radius $R$ is completely filled with water. Now an orifice of small area $a$ is made at the bottom of tank. (Given that $R=10$ m, $a=1\,\mathrm{mm^2}$)
\qfig{p1_phy_s3_para1_2_1}
\qn{1} Time rate of volume emission from orifice when tank have $\dfrac{h}{2}$ height of water level is (in $\mathrm{cm^3/sec}$)
\qn{2} The time required to empty the tank is $n\times10^6$ sec. Here $n$ is $-$
\paraq{3}{4}
A small object of mass 2 kg is placed on the x-axis in a potential field which potential energy is define as $U=16\,(x^2-2x)$ where position on x-axis, $x$ is in meter and potential energy $U$ is in joule. If the object is released from origin i.e. $x=0$. Then answer the following questions.
\qn{3} Angular frequency of oscillations ( in rad/second) is :
\qn{4} Maximum kinetic energy (in joule) during oscillations is :
\paraq{5}{6}
A block $A$ of mass 1 kg connected with a spring is executing SHM. The displacement time equation of block is $x=8+4\sin(10t)$ (here $x$ is in m). An identical block $B$ moving towards negative x-axis with velocity 10 m/s collides with block $A$ and sticks to it at time $t=0$
\qfig{p1_phy_s3_para5_6_1}
\qn{5} Speed of block $A$ (in m/s) just after the collision is
\qn{6} New amplitude (in m) of oscillation of block $A$ is
\sect{SECTION-II(ii)}
\qn{7} A spring of force constant $k=300$ N/m connects two blocks having masses 2 kg and 3 kg, lying on a smooth horizontal plane. If the spring block system is released from a stretched position, the number of complete oscillations in 1 minute is $6\times N$, find the value of $N$. Take $\pi=\sqrt{10}$.
\qn{8} A block of mass $m$ hangs from a vertical spring of spring constant $k$. If the block is pulled down by a distance of $3\,mg/k$ from its equilibrium position and released, for the subsequent SHM of the block the ratio of maximum elongation in the spring to maximum compression in it is
\qn{9} If natural frequency $\omega$ of the system shown in figure is $x\sqrt{\dfrac{k}{5m}}$, find $x$.
\qfig{p1_phy_s4_q09_1}
\partsec{2}{CHEMISTRY}
\sect{SECTION-I(i)}
\qn{1} Match the types of defect given in Column-I with the statement given in Column-II.
\par\nobreak
{\small\setlength{\tabcolsep}{3pt}\begin{center}
\begin{tabular}{|c|L{1.9cm}|c|L{3.1cm}|}
\hline
\multicolumn{2}{|c|}{\textbf{Column-I}} & \multicolumn{2}{c|}{\textbf{Column-II}}\\
\hline
(i) & Impurity defect & (a) & NaCl with anionic sites called F-centres\\
\hline
(ii) & Metal excess defect & (b) & FeO doped with Fe$^{3+}$\\
\hline
(iii) & Metal deficiency defect & (c) & NaCl doped with Sr$^{2+}$ and some cationic sites vacant\\
\hline
\end{tabular}\end{center}}
\begin{qoptions}
\opt{A}{(i)-a, (ii)-b, (iii)-c}
\opt{B}{(i)-c, (ii)-a, (iii)-b}
\opt{C}{(i)-a, (ii)-c, (iii)-b}
\opt{D}{(i)-b, (ii)-c, (iii)-a}
\end{qoptions}
\qn{2} Which of the following structure is non-planar and have atleast one `Cl' as bridging atom?
\begin{qoptions}
\opt{A}{$\mathrm{Be_2Cl_4}$}
\opt{B}{$\mathrm{Al_2Cl_6}$}
\opt{C}{$\mathrm{Si_2Cl_6}$}
\opt{D}{$\mathrm{I_2Cl_6}$}
\end{qoptions}
\qn{3} Which of following is optically active in nature
\begin{qoptions}
\opt{A}{\qfigi{p1_chem_s1_q03_A}}
\opt{B}{\qfigi{p1_chem_s1_q03_B}}
\opt{C}{\qfigi{p1_chem_s1_q03_C}}
\opt{D}{\qfigi{p1_chem_s1_q03_D}}
\end{qoptions}
\qn{4} Which one is correct in a given isomeric pair :
\begin{qoptions}
\opt{A}{\qfigi{p1_chem_s1_q04_A}}
\opt{B}{\qfigi{p1_chem_s1_q04_B}}
\opt{C}{\qfigi{p1_chem_s1_q04_C}}
\opt{D}{\qfigi{p1_chem_s1_q04_D}}
\end{qoptions}
\sect{SECTION-I(ii)}
\qn{5} A cubical crystal is such that the corners atoms are occupied by P atoms, face centers are occupied by Q atoms, alternate tetrahedral voids are occupied by R atoms and edge centers are occupied by S atoms. Identify the correct option(s) :
\begin{qoptions}
\opt{A}{If there are no defects in the crystal, then the formula should be $\mathrm{PQ_3S_4R_3}$}
\opt{B}{If all the atoms lying on one of the body diagonal are missing then formula should be $\mathrm{PQ_4,\ R_{8/3}S_4}$}
\opt{C}{If all the atoms lying on a line passing through edge centers and face-centre of a face are missing, then formula should be $\mathrm{P_2Q_5R_8S_5}$}
\opt{D}{If all the atoms lying on a line passing through body center and connecting opposite face centers are missing, then formula will be $\mathrm{PQ_2R_4S_3}$}
\end{qoptions}
\qn{6} Which of the following statements are correct -
\begin{qoptions}
\opt{A}{Amorphous solids are also called pseudo solids or supercooled liquids.}
\opt{B}{If a solid is very hard, and it is electrical insulator in solid as well as in molten state and melts at extremely high temperature, then it can be covalent or network solid}
\opt{C}{In Rhombohedral crystal system, only primitive unit cell is possible}
\opt{D}{In fluorite structure coordination number of cation is 4.}
\end{qoptions}
\qn{7} Consider four configurations of $\mathrm{B_2}$ and select correct statements(s)-\\
(I) $\sigma1s^2,\sigma^*1s^2,\sigma2s^2\ ,\ \sigma^*2s^2\ ,\sigma2p_x^2$\\
(II) $\sigma1s^2,\sigma^*1s^2,\ \sigma2s^2\ ,\ \sigma^*2s^2\ ,\pi2p_y^2$\\
(III) $\sigma1s^2\ ,\sigma^*1s^2,\sigma2s^2,\sigma^*2s^2,\pi2p_x^1=\pi2p_y^1$\\
(IV) $\sigma1s^3\ ,\sigma^*1s^3,\sigma2s^3,\sigma^*2s^1$
\begin{qoptions}
\opt{A}{(I) is ground state configuration}
\opt{B}{(II) violate Hund's rule}
\opt{C}{(III) violate s-p mixing}
\opt{D}{(IV) violate Pauli exclusion principle}
\end{qoptions}
\qn{8} Which of the following is/are nucleophile
\begin{qoptions}
\opt{A}{$:\mathrm{CCl_2}$}
\opt{B}{$\overset{\bullet\bullet}{\mathrm{N}}\mathrm{H_2}$}
\opt{C}{$\mathrm{CH_3{-}CH{=}CH_2}$}
\opt{D}{\qfigi{p1_chem_s2_q08_D}}
\end{qoptions}
\qn{9} Which of the following reaction(s) gives aromatic product as major product :
\begin{qoptions}
\opt{A}{\qfigi{p1_chem_s2_q09_A}}
\opt{B}{\qfigi{p1_chem_s2_q09_B}}
\opt{C}{\qfigi{p1_chem_s2_q09_C}}
\opt{D}{\qfigi{p1_chem_s2_q09_D}}
\end{qoptions}
\qfullbegin
\qn{10} Consider following molecular orbitals (MO) which are formed by LCAO of 2p-2p orbitals and select correct option
\qfig{p1_chem_s2_q10_1}
\qfullsplit
\begin{qoptions}
\opt{A}{MO-IV is a antibonding molecular orbital}
\opt{B}{In $\mathrm{F_2}$ MO-IV is the HOMO \& MO-II is the LUMO}
\opt{C}{In case of $\mathrm{O_2}$ one unpaired electron in MO-III and other one in MO-IV respectively}
\opt{D}{Energy of molecular orbitals in $\mathrm{O_2^{2-}}$ is I $<$ III $<$ IV $<$ II}
\end{qoptions}
\qfullend
\sect{SECTION-II(i)}
\paraq{1}{2}
A hypothetical metal M (Atomic weight $=450$) crystallizes as a cubic crystal such that metal atoms touch each other along body diagonal of cube. Edge length of cubic unit cell is 1nm.\\
If this metal is melted then it's density in molten form is 1.2 gm/cc. ($N_A=6\times10^{23}$)
\qn{1} Density of solid metal in gm/cc is -
\qn{2} Percentage of empty space in molten state of metal will be - \ [Use : $\pi\sqrt3=5.44$]
\paraq{3}{4}
In the heavier elements of p-block elements, however as we expect form the ``inert pair effect, the lower oxidation state becomes more common on descending the group.'' Covalent nature in ionic compound is explained on the basis of polarisation effect.
\qn{3} Atomic number of metal in the compound having most covalent character among the $\mathrm{ZnCl_2}$ , KCl, $\mathrm{CaCl_2}$ and NaCl?
\qn{4} Among the $\mathrm{Tl^{+x}}$, $\mathrm{Pb^{+x}}$ , $\mathrm{Bi^{+x}}$ , $\mathrm{Sn^{+x}}$ the maximum possible value of $x$ will be.\\
($x$ is most stable oxidation state of that element)
\paraq{5}{6}
Dehydration of alcohol is an elimination reaction, which may proceed via formation of carbocation intermediate \& carbocation rearrangement is also considered if required. Alkene is the product of the reaction where usually more stable alkene is major product.
\qfig{p1_chem_s3_para5_6_1}
\qn{5} Number of $\alpha$-H present in major product :
\qn{6} Number of T.S. (transition state) in given reaction for major product :
\sect{SECTION-II(ii)}
\qn{7} A crystalline solid of a pure substance has a face-centred cubic structure with a cell edge of 200 pm. If the density of the substance in the crystal is 8g $\mathrm{cm^{-3}}$, then the number of atoms present in 32g of the crystal is $N\times10^{24}$. The value of $N$ is
\qn{8} If radius of $\mathrm{M^+}$ is 113 pm and radius of $\mathrm{X^-}$ is 135 pm then the expected value of coordination number of $\mathrm{X^-}$ in MX(s) is ?
\qn{9} Which of the following pair having correct order of covalent character\\
($\mathrm{AlCl_3>MgCl_2}$), ($\mathrm{BeCl_2>LiCl}$), ($\mathrm{PbCl_4>PbCl_2}$), (HgS $>$ ZnS), ($\mathrm{NiCl_2>CaCl_2}$)
\partsec{3}{MATHEMATICS}
\sect{SECTION-I(i)}
\qn{1} Let $f(x)=e^x+e^{-x}+a$ has exactly 2 distinct real roots, then
\begin{qoptions}
\opt{A}{$a>0$}
\opt{B}{$a<0$}
\opt{C}{$a<-2$}
\opt{D}{$a\in\phi$}
\end{qoptions}
\qn{2} If product of roots `$p$' of equation $x^2-2ax+8-a=0$ lies between its roots, then maximum integral value of `$p$' is -
\begin{qoptions}
\opt{A}{1}
\opt{B}{4}
\opt{C}{6}
\opt{D}{7}
\end{qoptions}
\qn{3} If $x$ is an odd integer and $y,z$ are even integers, satisfying $x+y+z=47$ and $x,y,z\ge0$, then number of solutions are
\begin{qoptions}
\opt{A}{100}
\opt{B}{200}
\opt{C}{300}
\opt{D}{400}
\end{qoptions}
\qn{4} The term independent of $x$ in the expansion of $\left(x^2+x+1+\dfrac1x\right)^5$ is equal to -
\begin{qoptions}
\opt{A}{98}
\opt{B}{99}
\opt{C}{100}
\opt{D}{101}
\end{qoptions}
\sect{SECTION-I(ii)}
\qn{5} If the $9^{th}$ term in the expansion of $\left(\dfrac{1}{x^2}+\dfrac{x}{2}\log_2x\right)^{12}$ is equal to 495, then the value of $x$ can be-
\begin{qoptions}
\opt{A}{1}
\opt{B}{An integer greater than 1}
\opt{C}{a fraction}
\opt{D}{an irrational number}
\end{qoptions}
\qn{6} If $f(x)=2x^2+kx+50$, $k\ge20$ and magnitude of difference of the roots of the equation $f(x)=0$ is 8, then which of the following can be true -
\begin{qoptions}
\opt{A}{minimum value of $f(x)$ is $-32$}
\opt{B}{minimum value of $f(x)$ is $-64$}
\opt{C}{graph of $f(x)$ cuts positive y-axis}
\opt{D}{graph of $f(x)$ cuts negative y-axis}
\end{qoptions}
\qn{7} Sum of the series $({}^{100}C_1)^2+2.({}^{100}C_2)^2+3.({}^{100}C_3)^2+\dots+100({}^{100}C_{100})^2$ equals-
\begin{qoptions}
\opt{A}{$\dfrac{2^{99}\,[1.3.5.....(199)]}{99!}$}
\opt{B}{$100.{}^{200}C_{100}$}
\opt{C}{$50.{}^{200}C_{100}$}
\opt{D}{$100.{}^{199}C_{99}$}
\end{qoptions}
\qn{8} Let
\[f(x)=\begin{cases}x^2-3x+4\ ;&x<3\\ x+7\ ;&x\ge3\end{cases}\]
and
\[g(x)=\begin{cases}x+6\ ;&x<4\\ x^2+x+2\ ;&x\ge4\end{cases}\]
then which of the following is/are true -
\begin{qoptions}
\opt{A}{$(f+g)(1)=9$}
\opt{B}{$(f-g)(3.5)=1$}
\opt{C}{$(fg)(0)=24$}
\opt{D}{$\left(\dfrac fg\right)(5)=\dfrac83$}
\end{qoptions}
\qn{9} If $a,b,c$ are three consecutive even natural number, then which of the following is (are) true -
\begin{qoptions}
\opt{A}{$\log_{10}(4+ac)=2\log_{10}b$}
\opt{B}{$\log_{10}(a+c)=\log_{10}2+\log_{10}b$}
\opt{C}{$\log_{10}(a^2+c^2)=3\log_{10}2+\log_{10}b$}
\opt{D}{$\log_{10}(a+b+c)=\log_{10}6+\log_{10}b$}
\end{qoptions}
\qn{10} If $\underline{|500}=k.\left(2^a3^b5^c\right)$, (where $a,b,c\in N$, $k$ does not contain any factor of 2,3 and 5) then-
\begin{qoptions}
\opt{A}{$a>c$}
\opt{B}{$c=124$}
\opt{C}{number of zeros at the end of $\underline{|500}$ is equal to 124}
\opt{D}{$\underline{|500}$ is an even number.}
\end{qoptions}
\sect{SECTION-II(i)}
\paraq{1}{2}
Consider the name of a baby girl `HARSHITA'. New names are formed (with or without meaning) by using all the letters of the given word.
\qn{1} Number of words in which vowels occupy even places -
\qn{2} Number of words which contains the word `AAI' is-
\paraq{3}{4}
Let `$A$' is the number of solution(s) of the equation
\[x^{\log x\,\log 4}\times4^{\log x^3}+\dfrac{64}{4^{\log^2x}\times x^{3\log4}}=65,\]
$B$ is the value of $\sec^2x$ if $\dfrac{\sin^2x-2\cos^2x+1}{\sin^2x+2\cos^2x-1}=4$ and $C$ is the value of the expression $\cos2x-\dfrac{\sin x\cos3x}{\sin2x}$, (wherever defined).
\qn{3} The value of $(B+C)$ is -
\qn{4} Least integral value of $(Ax^2+Bx+C)$ is -
\paraq{5}{6}
Let $(1-x+x^2)^{30}=a_0+a_1x+a_2x^2+a_3x^3+\dots\dots\dots+a_{60}x^{60}$.
\qn{5} $\left(\dfrac{a_0^2-a_1^2+a_2^2-a_3^2+\dots\dots\dots}{a_{30}}\right)$
\qn{6} $4\left(\dfrac{a_0+a_2+a_4+\dots\dots\dots}{(3^{30}+1)}\right)$
\sect{SECTION-II(ii)}
\qn{7} Let sum of binomial coefficient in the expansion of $(x+2y)^n$ is 1024 ($n\in N$), then number of divisors of $(120)(n)$ divisible by 75 is
\qn{8} If $\alpha,\beta,\gamma$ are the roots of the equation $x^3-5x^2+3=0$, then the value of $(\alpha-2)^3(\beta-2)^3(\gamma-2)^3$ is
\qn{9} If L.C.M. of $(a,b)=5^2\times3^4$, L.C.M. of $(b,c)=5^3\times3^4$ \& L.C.M. of $(c,a)=5^3\times3^4$, then possible ordered triplets of $(a,b,c)$ are $k$ then $(k-58)$ is equal to
\newpage
\papersec{2}
\partsec{1}{PHYSICS}
\sect{SECTION-I}
\qn{1} A tray of the shape is shown in figure. It is completely filled with a liquid of density $\rho$. Find the force on its rectangular wall $ABCD$ due to liquid. The wall is inclined at an angle $\theta$ to the vertical. (Ignore atmospheric pressure)
\qfig{p2_phy_s1_q01_1}
\begin{qoptions}
\opt{A}{$\dfrac12\rho gh^2b$}
\opt{B}{$\dfrac12\rho gh^2b\sec\theta$}
\opt{C}{$\rho gh^2b$}
\opt{D}{$\dfrac12\rho gh^2b\tan\theta$}
\end{qoptions}
\qn{2} A block of mass $m$ is restricted to move in a smooth groove. It is connected to a spring of spring constant $k$. Neglecting the effect of gravity find time period of small oscillation :-
\qfig{p2_phy_s1_q02_1}
\begin{qoptions}
\opt{A}{$\dfrac{2\pi}{\cos\theta}\sqrt{\dfrac mk}$}
\opt{B}{$2\pi\sqrt{\dfrac mk}$}
\opt{C}{$\dfrac{2\pi}{\sin\theta}\sqrt{\dfrac mk}$}
\opt{D}{$2\pi\sqrt{\dfrac{m\cos^2\theta}{k}}$}
\end{qoptions}
\qn{3} The time period of a particle in simple harmonic motion is equal to the time between consecutive appearances of the particle at a particular point in its motion. This point is
\begin{qoptions}
\opt{A}{The mean position}
\opt{B}{Between the mean position and the positive extreme}
\opt{C}{An extreme position}
\opt{D}{Between the mean position and the negative extreme.}
\end{qoptions}
\qn{4} A thin tube in vertical plane is shown in figure. It is filled with a liquid of density $\rho$ and its end $B$ is closed. Then the force exerted by the liquid on the tube at end $B$ will be : [Neglect atmospheric pressure and the radius of the tube to be negligible in comparison to $\ell$]
\qfig{p2_phy_s1_q04_1}
\begin{qoptions}
\opt{A}{$\rho.g.\ell.A$}
\opt{B}{$3\rho g\ell A$}
\opt{C}{$\dfrac{\rho g\ell A}{2}$}
\opt{D}{Zero}
\end{qoptions}
\qn{5} Potential energy associated with a particle of mass $m$ moving along $x$ axis is given by $U=U_0[1-\cos(bx)]$ where $U_0$ \& $b$ are positive constant. Find the time period of small oscillation
\begin{qoptions}
\opt{A}{$2\pi\sqrt{\dfrac{m}{U_0b}}$}
\opt{B}{$2\pi\sqrt{\dfrac{mb}{U_0}}$}
\opt{C}{$2\pi\sqrt{\dfrac{m}{U_0}}$}
\opt{D}{$\dfrac{2\pi}{b}\sqrt{\dfrac{m}{U_0}}$}
\end{qoptions}
\qn{6} The total mechanical energy of a spring-mass system in simple harmonic motion is $E=\dfrac12m\omega^2A^2$. Suppose the oscillating particle is replaced by another particle of double the mass while the amplitude $A$ remains the same. The new mechanical energy will
\begin{qoptions}
\opt{A}{become $2E$}
\opt{B}{become $E/2$}
\opt{C}{become $\sqrt2E$}
\opt{D}{remain $E$}
\end{qoptions}
\qn{7} An L shaped thin tube contains a liquid of density $\rho$. It is moved horizontally with a constant acceleration $\dfrac g4$. Part $AB$ of tube is vertical and $BC$ is horizontal. End $A$ is open to atmosphere and atmospheric pressure is $P_0$. If pressure of liquid at point $C$ (just inside the tube) is equal to $P_0$ then:-
\qfig{p2_phy_s1_q07_1}
\begin{qoptions}
\opt{A}{$L=4H$}
\opt{B}{$L=2H$}
\opt{C}{$L=3H$}
\opt{D}{$L=H$}
\end{qoptions}
\qn{8} In the figure block lies on smooth horizontal surface and both the springs have same spring constant $k$. The spring $B$ is attached with block while spring $A$ just touches the block as the position shown and both the springs are in relaxed state. The time period of oscillation is :-
\qfig{p2_phy_s1_q08_A}
\begin{qoptions}
\opt{A}{$2\pi\sqrt{\dfrac{m}{2k}}$}
\opt{B}{$2\pi\sqrt{\dfrac{3m}{2k}}$}
\opt{C}{$2\pi\sqrt{\dfrac mk}$}
\opt{D}{$\dfrac{(\sqrt2+1)}{\sqrt2}\pi\sqrt{\dfrac mk}$}
\end{qoptions}
\qn{9} A particle executes SHM about the point $x=0$. At time $t=0$ it has position $x=A$ and zero velocity. The angular frequency of oscillation is $\omega$. The acceleration of the particle at $t=\dfrac{\pi}{4\omega}$ is :-
\begin{qoptions}
\opt{A}{$-A\omega^2$}
\opt{B}{$-\dfrac{A\omega^2}{\sqrt2}$}
\opt{C}{Zero}
\opt{D}{$-A\omega^2\sqrt2$}
\end{qoptions}
\qn{10} Two springs each of spring constant $K$ are connected with the centre of a sphere of mass $m$ and radius $R$ as shown in figure. Given system lies on horizontal surface and both the springs are horizontal. Friction is sufficient to prevent relative slipping between horizontal surface and sphere. From the position where both springs are at natural length, sphere is slightly displaced by a distance `$x_0$' and released from rest. After being released from rest find minimum time taken by sphere to move by a distance $\dfrac{x_0}{2}$.
\qfig{p2_phy_s1_q10_1}
\begin{qoptions}
\opt{A}{$\dfrac\pi3\sqrt{\dfrac{7m}{10k}}$}
\opt{B}{$\dfrac\pi3\sqrt{\dfrac{5m}{7k}}$}
\opt{C}{$\dfrac\pi4\sqrt{\dfrac{7m}{5k}}$}
\opt{D}{$\dfrac\pi6\sqrt{\dfrac{5m}{7k}}$}
\end{qoptions}
\qn{11} A simple pendulum is suspended from the ceiling of a car moving with constant speed on a horizontal circular track of radius $R$. Speed of the car is given by $v=\sqrt{gR}$ where $g$ is acceleration due to gravity and $R$ is radius of circular turn. Find the time period of a simple pendulum of length $\ell$ suspended from the ceiling of the car.
\begin{qoptions}
\opt{A}{$2\pi\sqrt{\dfrac{\ell}{2g}}$}
\opt{B}{$2\pi\sqrt{\dfrac{2\ell}{g}}$}
\opt{C}{$2\pi\sqrt{\dfrac{\ell}{g}}$}
\opt{D}{$2\pi\sqrt{\dfrac{\ell}{\sqrt2g}}$}
\end{qoptions}
\qn{12} In a horizontal spring-mass system, at $t=0$, mass `$m$' is released from rest after being displaced (from NLP of spring) towards right by some distance on a frictionless surface. Find minimum time after which speed of block becomes $\dfrac{1}{\sqrt2}$ times of maximum possible speed during the motion.
\qfig{p2_phy_s1_q12_A}
\begin{qoptions}
\opt{A}{$\dfrac\pi2\sqrt{\dfrac mk}$}
\opt{B}{$\dfrac\pi4\sqrt{\dfrac mk}$}
\opt{C}{$\pi\sqrt{\dfrac mk}$}
\opt{D}{$\dfrac\pi8\sqrt{\dfrac mk}$}
\end{qoptions}
\qn{13} Two particles $A$ and $B$ are performing SHM along same straight line with same amplitude $a$ and same angular frequency $\omega$. At $t=0$, $A$ is at the left extreme and $B$ is at a distance $a/2$ from mean position to the right of mean position and moving toward mean position. Then:
\begin{qoptions}
\opt{A}{their phase difference is $\dfrac{5\pi}{3}$}
\opt{B}{their phase difference is $\dfrac{\pi}{3}$}
\opt{C}{time after which they cross each other is $\dfrac{2\pi}{3\omega}$}
\opt{D}{time after which they cross each other is $\dfrac{\pi}{3\omega}$}
\end{qoptions}
\qn{14} Figure shows the potential-energy diagram and the total energy line of a particle oscillating on a spring. Which of the following option is incorrect.
\qfig{p2_phy_s1_q14_1}
\begin{qoptions}
\opt{A}{the amplitude of oscillation of particle is 6 cm.}
\opt{B}{the spring constant is 4.44 kN/m.}
\opt{C}{the mass of particle is 1kg.}
\opt{D}{the maximum kinetic energy of particle is at $x=20$ cm.}
\end{qoptions}
\qn{15} The coefficient of friction between block of mass $m$ and $2m$ is $\mu=2\tan\theta$. There is no friction between block of mass $2m$ and inclined plane. The maximum amplitude of two block system for which there is no relative motion between both the blocks is:
\qfig{p2_phy_s1_q15_1}
\begin{qoptions}
\opt{A}{$g\sin\theta\sqrt{\dfrac km}$}
\opt{B}{$\dfrac{mg\sin\theta}{k}$}
\opt{C}{$\dfrac{3mg\sin\theta}{k}$}
\opt{D}{None of these}
\end{qoptions}
\qn{16} A block of ice is floating in a liquid (in beaker) of specific gravity 1.2, which is immiscible with water. When all the ice melts, the liquid level in the beaker.
\begin{qoptions}
\opt{A}{falls}
\opt{B}{rises}
\opt{C}{remains unchanged}
\opt{D}{None}
\end{qoptions}
\qn{17} A rod of mass `$m$' and length `$L$' is suspended in vertical plane so as to oscillate freely about an axis normal to the plane of paper passing through `$S$' (point of suspension). Find the distance of `$S$' from end $A$ of rod such that time period of small oscillation is minimum.
\qfig{p2_phy_s1_q17_1}
\begin{qoptions}
\opt{A}{$\dfrac L2$}
\opt{B}{$\dfrac{L}{\sqrt{12}}$}
\opt{C}{$\dfrac L2+\dfrac{L}{\sqrt{12}}$}
\opt{D}{$\dfrac L2-\dfrac{L}{\sqrt{12}}$}
\end{qoptions}
\qn{18} A wooden cube (density of wood `$\dfrac\rho3$') of side `$\ell$' floats in a liquid of density `$\rho$' with its upper and lower surfaces horizontal. If the cube is pushed slightly down and released, it performs simple harmonic motion of period `$T$'. Then, `$T$' is equal to :-
\begin{qoptions}
\opt{A}{$2\pi\sqrt{\dfrac{3\ell}{2g}}$}
\opt{B}{$2\pi\sqrt{\dfrac{\ell}{3g}}$}
\opt{C}{$2\pi\sqrt{\dfrac{3\ell}{g}}$}
\opt{D}{$2\pi\sqrt{\dfrac{2\ell}{g}}$}
\end{qoptions}
\qn{19} A small block of mass $m$ is kept on a bigger block of mass $M$ which is attached to a vertical spring of spring constant $k$ as shown in the figure. The system oscillates vertically. What can be the maximum amplitude with which the two blocks may oscillate together ?
\qfig{p2_phy_s1_q19_1}
\begin{qoptions}
\opt{A}{$\dfrac{2(M+m)g}{k}$}
\opt{B}{$\dfrac{(m)g}{k}$}
\opt{C}{$\dfrac{(M)g}{k}$}
\opt{D}{$\dfrac{(M+m)g}{k}$}
\end{qoptions}
\qn{20} A block is floating in water in a container at rest. Length of block inside the container is $h$ as shown in figure. Density of block is $\dfrac\rho4$ and density of water is $\rho$. Now the container starts accelerating upwards with constant acceleration $\dfrac g4$. The new length of block inside the water will be [consider container achieves acceleration slowly]
\qfig{p2_phy_s1_q20_1}
\begin{qoptions}
\opt{A}{$\dfrac h4$}
\opt{B}{$\dfrac h2$}
\opt{C}{$h$}
\opt{D}{$2h$}
\end{qoptions}
\sect{SECTION-II}
\qn{1} A solid of density $\dfrac\rho2$ ($\rho$ = density of water) is held inside a water tank, tied to a string, as shown in figure. The tank is accelerated uniformly in horizontal direction such that the water surface makes an angle $45^\circ$ with the horizontal. Find the angle $\theta$ (in degree) that the string makes with base of container in equilibrium.
\qfig{p2_phy_s2_q01_1}
\qn{2} Liquid is flowing in a horizontal pipe. Two tubes, $A$ \& $B$ are inserted in the pipe. Tube $B$ is L-shaped and its cross-section at 1 is held normal to the flow direction. Liquid rises in both tubes. Height of liquid column in $B$ is $h$ more than that in $A$. Considering $h=5$ cm and $g=10\ \mathrm{m/s^2}$, find the speed of liquid flow ($v$) in m/sec in the pipe.
\qfig{p2_phy_s2_q02_1}
\qn{3} Equation of motion of a particle performing SHM is given by\\
$y=\dfrac{10}{\pi}\sin\left[2000\pi t-\dfrac\pi6\right]$, where $y$ is in cm and $t$ is in sec.\\
If time period of SHM is $t$ and maximum possible velocity of the particle during SHM is $v_{\max}$ then find the value of $t\times v_{\max}$ (in SI unit).
\qn{4} A uniform disc of radius 5.0 cm and mass 200g is fixed at its centre to a metal wire, the other end of which is fixed with a clamp. The hanging disc is rotated about the wire through an angle and a released. If the disc makes torsional oscillations with time period 0.20 s, find the torsional constant (in S.I. units) of the wire. (Take : $\pi^2=10$)
\qfig{p2_phy_s2_q04_1}
\qn{5} A light wire $AB$ of length 10 cm can slide on a vertical frame as shown in figure. There is a film of soap solution trapped between the frame and the wire. Find the load $W$ (in g) that should be suspended from the wire to keep it in equilibrium. Neglect friction. Surface tension of soap solution $=25$ dyne/cm. Take $g=10\ \mathrm{m/s^2}$.
\qfig{p2_phy_s2_q05_1}
\qn{6} There is an air bubble of radius 1.0 mm in a liquid of surface tension 0.075 N/m and density $1000\ \mathrm{kg/m^3}$. The bubble is at a depth of 10 cm below the free surface. By what amount is the pressure (in $\mathrm{N/m^2}$) inside the bubble greater than the atmospheric pressure ? Take : $g=9.8\ \mathrm{m/s^2}$.
\qn{7} A narrow capillary tube of radius 0.5 mm contains some water in it. When tube is held vertical, the water remains inside the tube as shown. Find the radius of curvature (in mm) of the lower water surface. (Surface tension $s=0.075$ N/m, contact angle $=0^\circ$, $g=10\ \mathrm{m/s^2}$).
\qfig{p2_phy_s2_q07_1}
\qn{8} A simple pendulum oscillating in air has period $T$. The bob of the pendulum is completely immersed in a non-viscous liquid. The density of the liquid is $\dfrac1{16}$ th of the density of material of the bob. If bob is inside liquid all the time and its period of oscillation in this liquid is given by $4T\sqrt{\dfrac1N}$ then find the value of $N$.
\qn{9} A block $A$ is placed on a table. There is a 3.0 mm thick layer of oil between the table and block $A$. Face of the block in contact with oil has area $0.1\ \mathrm{m^2}$. The block $A$ is connected to another block $B$ of mass $m=0.01$ kg by massless string passing over a smooth pulley. The system is released from rest. It accelerates and finally attains a constant speed of 0.085 m/s. Find coefficient of viscosity of liquid (in S.I. Units). Neglect viscous force and buoyancy force due to air on the blocks ($g=9.8\ \mathrm{m/s^2}$)
\qfig{p2_phy_s2_q09_1}
\qn{10} Two small balls each of mass $m$ are connected by a uniform rod of mass $m$ and length $L$. The given system is suspended from its centre by a thin wire of torsional constant $k$. The rod is rotated about the wire through an angle $\theta_0$ and released. Find the velocity of each ball as the system passes through the mean position. If this velocity is given by $v=\theta_0\sqrt{\dfrac{\alpha k}{7m}}$ then fill the value of $\alpha$.
\qfig{p2_phy_s2_q10_1}
\partsec{2}{CHEMISTRY}
\sect{SECTION-I}
\qn{1} If the value of ionic radius ratio $\left(\dfrac{r^+}{r^-}\right)$ is 0.52 in an ionic compound, the geometrical arrangement of ions in crystal is
\begin{qoptions}
\opt{A}{Tetrahedral}
\opt{B}{Planar}
\opt{C}{Octahedral}
\opt{D}{Pyramidal}
\end{qoptions}
\qn{2} Identical 4 spheres are taken and are arranged in a layer of square packing touching each other as shown
\qfig{p2_chem_s1_q02_1}
The percentage of vacant space is (in 3-D)
\begin{qoptions}
\opt{A}{$100\left(1-\dfrac{3\pi}{8}\right)$}
\opt{B}{$100\left(1-\dfrac\pi6\right)$}
\opt{C}{$100-\dfrac{3\pi}{8}$}
\opt{D}{$\dfrac\pi6$}
\end{qoptions}
\qn{3} If $r_{\mathrm{Na^+}}=95$ pm and $r_{\mathrm{Cl^-}}=181$ pm in NaCl (rock salt) structure. What is the shortest distance between $\mathrm{Na^+}$ ions ?
\begin{qoptions}
\opt{A}{778 pm}
\opt{B}{276 pm}
\opt{C}{195 pm}
\opt{D}{390 pm}
\end{qoptions}
\qn{4} In a hypothetical solid `C' atoms are found to form cubical close packed lattice. `A' atoms occupy all tetrahedral voids \& `B' atoms occupy all octahedral voids. `A' and `B' atoms
\qfig{p2_chem_s1_q04_1}
are of appropriate size i.e. there is no distortion in ccp lattice of C atoms. Now if a plane as shown in the following figure is cut, then the cross section of this plane will look like.
\begin{qoptions}
\opt{A}{\qfigi{p2_chem_s1_q04_A}}
\opt{B}{\qfigi{p2_chem_s1_q04_B}}
\opt{C}{\qfigi{p2_chem_s1_q04_C}}
\opt{D}{\qfigi{p2_chem_s1_q04_D}}
\end{qoptions}
\qn{5} 120 g glucose $[\mathrm{C_6H_{12}O_6(s)}]$ is dissolved to prepare 1 L aqueous solution of density 1.2g/mL. Identify incorrect option for the solution
\begin{qoptions}
\opt{A}{Molarity $=2/3$ M}
\opt{B}{\% w/w $=10\%$}
\opt{C}{\% w/v $=12\%$}
\opt{D}{Molality $=1.85$ mol/kg}
\end{qoptions}
\qn{6} In an ionic solid atoms `A' form hexagonal close packing crystal with atoms `B' in ($2/3^{\mathrm{rd}}$) octahedral voids. Empirical formula of solid is -
\begin{qoptions}
\opt{A}{AB}
\opt{B}{$\mathrm{A_3B_2}$}
\opt{C}{$\mathrm{A_2B}$}
\opt{D}{$\mathrm{A_2B_3}$}
\end{qoptions}
\qn{7} p$\pi$ - d$\pi$ bond is present in :
\begin{qoptions}
\opt{A}{$\mathrm{BF_2Cl}$}
\opt{B}{$\mathrm{SiF_4}$}
\opt{C}{CO}
\opt{D}{$\mathrm{BeCl_2}$}
\end{qoptions}
\qn{8} In which of the following species odd electron present in pure orbital?
\begin{qoptions}
\opt{A}{$\mathrm{NO_2}$}
\opt{B}{$\mathrm{ClO_2}$}
\opt{C}{$\mathrm{CF_3}$}
\opt{D}{$\mathrm{ClO_3}$}
\end{qoptions}
\qn{9} In which of the following processes, the value of magnetic moment does not change ?
\begin{qoptions}
\opt{A}{$\mathrm{N_2\to N_2^-}$}
\opt{B}{$\mathrm{N_2\to N_2^+}$}
\opt{C}{$\mathrm{O_2\to O_2^+}$}
\opt{D}{$\mathrm{O_2^+\to O_2^-}$}
\end{qoptions}
\qn{10} The nature of overlap in the bridge bond present in $\mathrm{I_2Cl_6}$ is :
\begin{qoptions}
\opt{A}{$\mathrm{sp^3}$--p--$\mathrm{sp^3}$}
\opt{B}{$\mathrm{sp^3d^2}$--$\mathrm{sp^3}$--$\mathrm{sp^3d^2}$}
\opt{C}{$\mathrm{sp^3d}$--$\mathrm{sp^3}$--$\mathrm{sp^3d}$}
\opt{D}{$\mathrm{sp^3d}$--$\mathrm{sp^2}$--$\mathrm{sp^3d}$}
\end{qoptions}
\qn{11} Which of the following is not a correct order\\
(i) AgF $<$ AgCl $<$ AgBr (Colour intensity)\\
(ii) $\mathrm{Rb_2CO_3>K_2CO_3>Na_2CO_3}$ (decomposition temperature)\\
(iii) $\mathrm{BeC_2O_4<MgC_2O_4<CaC_2O_4}$ (Solubility)\\
(iv) $\mathrm{Li_3N>Na_3N>K_3N}$ (Thermal stability)
\begin{qoptions}
\opt{A}{i, iii, iv}
\opt{B}{i, iii}
\opt{C}{only iii}
\opt{D}{iii and iv}
\end{qoptions}
\qn{12} Which of the following order is correct for covalent character (if the size of cations are identical in pairs)-
\begin{qoptions}
\opt{A}{$\mathrm{NaCl<CuCl}$}
\opt{B}{$\mathrm{CaCl_2<PdCl_2}$}
\opt{C}{$\mathrm{RbCl<AuCl}$}
\opt{D}{All of these}
\end{qoptions}
\qn{13} Select \textbf{INCORRECT} order against property indicated--
\begin{qoptions}
\opt{A}{$\mathrm{AsCl_3<SbCl_3<BiCl_3}$ (relative stability)}
\opt{B}{$\mathrm{ZnO>MgO>CaO>BaO}$ (Covalent character)}
\opt{C}{$\mathrm{NaF<MgF_2<AlF_3}$ (M.P.) Melting point}
\opt{D}{$\mathrm{BeF_2>MgF_2>CaF_2>SrF_2>BaF_2}$ (water solubility)}
\end{qoptions}
\qn{14} Which of the following pair of compounds doesn't represent enantiomeric pair :
\begin{qoptions}
\opt{A}{\qfigi{p2_chem_s1_q14_A}}
\opt{B}{\qfigi{p2_chem_s1_q14_B}}
\opt{C}{\qfigi{p2_chem_s1_q14_C}}
\opt{D}{\qfigi{p2_chem_s1_q14_D}}
\end{qoptions}
\qn{15} \qfigi{p2_chem_s1_q15_1}\par
Identify total number of 1,2-shift for above reaction ?
\begin{qoptions}
\opt{A}{2}
\opt{B}{3}
\opt{C}{4}
\opt{D}{5}
\end{qoptions}
\qn{16} Which of the following compound(s) contains plane of symmetry as well as centre of symmetry
\begin{qoptions}
\opt{A}{$\mathrm{CH_2Cl_2}$}
\opt{B}{\qfigi{p2_chem_s1_q16_B}}
\opt{C}{Cis-but-2-ene}
\opt{D}{\qfigi{p2_chem_s1_q16_2}}
\end{qoptions}
\qn{17} Which of the following compound(s) can not exist in meso form ?
\begin{qoptions}
\opt{A}{\qfigi{p2_chem_s1_q17_A}}
\opt{B}{\qfigi{p2_chem_s1_q17_B}}
\opt{C}{\qfigi{p2_chem_s1_q17_C}}
\opt{D}{\qfigi{p2_chem_s1_q17_D}}
\end{qoptions}
\qn{18} The order of increasing nucleophilicity of the anions (in polar protic solvent)
\qfig{p2_chem_s1_q18_1}
\begin{qoptions}
\opt{A}{V $<$ I $<$ III $<$ II $<$ IV}
\opt{B}{III $<$ I $<$ IV $<$ II $<$ V}
\opt{C}{III $<$ II $<$ I $<$ IV $<$ V}
\opt{D}{II $<$ III $<$ I $<$ V $<$ IV}
\end{qoptions}
\qn{19} Which of the following order of leaving group is incorrect ?
\begin{qoptions}
\opt{A}{$\mathrm{F^-<Cl^-<Br^-<I^-}$}
\opt{B}{\qfigi{p2_chem_s1_q19_B}}
\opt{C}{$\mathrm{CH_3-OH>CH_3-O^-}$}
\opt{D}{$\mathrm{Cl^-<H^-}$}
\end{qoptions}
\qn{20} The most stable carbocation after rearrangement is
\qfig{p2_chem_s1_q20_1}
\begin{qoptions}
\opt{A}{\qfigi{p2_chem_s1_q20_A}}
\opt{B}{\qfigi{p2_chem_s1_q20_B}}
\opt{C}{\qfigi{p2_chem_s1_q20_C}}
\opt{D}{\qfigi{p2_chem_s1_q20_D}}
\end{qoptions}
\sect{SECTION-II}
\qn{1} 2M of 100 ml $\mathrm{Na_2SO_4}$ is mixed with 3M of 100 ml NaCl solution and 1M of 200 ml $\mathrm{CaCl_2}$ solution. Then find the ratio of the molarity of all cation and all anions in final solution ?
\qn{2} In an ionic solid $r_{(+)}=1.6$ \AA\ and $r_{(-)}=1.864$ \AA. Use the radius ratio rule to determine the edge length of the cubic unit cell in \AA. [ $\sqrt3=1.732$ ]
\qn{3} A sample of $\mathrm{H_2O_2}$ is labelled `28 volumes' at 1 atm \& 273K. Find molality of the solution if its density is 265 g/L ?
\qn{4} For any metal $^{18}$M crystallising in fcc arrangement (Atomic radius $=1/\sqrt2$ \AA) with density is 67.5 gm/cc hence atomic mass of the metal is [$N_A=6\times10^{23}$] (in gram/mole)
\qn{5} Write the number of pairs in which size of first element or ion is higher as compare to II$^{\mathrm{nd}}$ out of following pairs :\\
(O, S), (He, Ne), (Kr, Ne), ($\mathrm{Na}$, $\mathrm{Na^+}$), (Cl, $\mathrm{Cl^-}$), (Li, Na), ($\mathrm{Li^+}$, $\mathrm{Na^+}$)
\qn{6} Find the number of element having first ionization energy lower than Be.\\
Mg, Ca, Sr, Ba, Li
\qn{7} How many orders are CORRECT\\
(a) Acidic strength $\Rightarrow$ ClOH $<$ BrOH $<$ IOH\\
(b) Basic strength $\Rightarrow$ MgO $<$ CaO $<$ SrO $<$ BeO\\
(c) Electronegativity $\Rightarrow$ I $<$ Br $<$ N $<$ O $<$ F\\
(d) Electron affinity $\Rightarrow$ Mg $<$ Na $<$ Si $<$ S $<$ Cl\\
(e) \% Ionic character $\Rightarrow$ NaCl $<$ KF $<$ RbF
\qn{8} Among the following, the number of molecules which are optically active \& possess $C_2$ axis of symmetry are :
\qfig{p2_chem_s2_q08_1}
$\mathrm{CHCl_3}$ , 2,5-dimethylthiophene,
\qfig{p2_chem_s2_q08_2}
\qn{9} How many compound(s) can show geometrical isomerism ?
\qfig{p2_chem_s2_q09_1}
\qn{10} How many of the given carbocation can undergo rearrangement :
\qfig{p2_chem_s2_q10_1}
\partsec{3}{MATHEMATICS}
\sect{SECTION-I}
\qn{1} A triangle is divided into 4 pieces and the area of the 4 regions are given inside the corresponding regions, then value of $x$ is -
\qfig{p2_math_s1_q01_1}
\begin{qoptions}
\opt{A}{18}
\opt{B}{15}
\opt{C}{22}
\opt{D}{12.5}
\end{qoptions}
\qn{2} A point $P$ is taken outside $\Delta ABC$ where $(B(1,\sqrt3)$, $A(0,0)$ and $C(2,0))$, but inside the acute angle $BAC$, such that $\angle APC=\dfrac\pi6$ and $\angle BPA=\dfrac\pi{12}$. Slope of the line $BP$ is
\begin{qoptions}
\opt{A}{$\sqrt3$}
\opt{B}{$-\sqrt3$}
\opt{C}{$\dfrac1{\sqrt3}$}
\opt{D}{$\dfrac{-1}{\sqrt3}$}
\end{qoptions}
\qn{3} Sum of all possible integral solutions of the inequality $(x-2)^2-5|x-2|+6<0$ is -
\begin{qoptions}
\opt{A}{4}
\opt{B}{5}
\opt{C}{7}
\opt{D}{8}
\end{qoptions}
\qn{4} If $x,y,z$ are non-zero real numbers then the minimum value of expression
\[\left(\dfrac{(x^8+4x^4+1)(y^4+3y^2+1)(z^2+2z+2)}{x^4y^2}\right)\]
equals -
\begin{qoptions}
\opt{A}{12}
\opt{B}{24}
\opt{C}{30}
\opt{D}{48}
\end{qoptions}
\qn{5} A man moves outside a park $(ABCDE)$ in such a way that he always remains at a distance of 3m from the wall (as shown in figure), then total distance travelled for a complete round is -
\qfig{p2_math_s1_q05_1}
\begin{qoptions}
\opt{A}{$(200+6\pi)$ m}
\opt{B}{$(200+8\pi)$ m}
\opt{C}{$(200+10\pi)$ m}
\opt{D}{$(200+4\pi)$ m}
\end{qoptions}
\qn{6} Coefficient of $x^5$ in the expansion of
\[\left(x+\dfrac1x\right)^{10}\left(x-\dfrac1x\right)^{9}\]
is -
\begin{qoptions}
\opt{A}{${}^9C_3$}
\opt{B}{${}^9C_4$}
\opt{C}{$-{}^9C_3$}
\opt{D}{$-{}^9C_4$}
\end{qoptions}
\qn{7} If sum of the last 30 coefficients in the expansion of $(1+x)^{59}$, when expanded in ascending power of `$x$' is $2^n$ then number of divisors of `$n$' of the form `$4\lambda+2$' ($\lambda\in N$) is -
\begin{qoptions}
\opt{A}{1}
\opt{B}{0}
\opt{C}{2}
\opt{D}{4}
\end{qoptions}
\qn{8} If $(1-x^2+x^3)(1+x)^n=a_0+a_1x+a_2x^2+\dots$, $a_i$'s $\in R$, then the value of $a_1^2-a_1-2a_2$ is -
\begin{qoptions}
\opt{A}{$-2$}
\opt{B}{2}
\opt{C}{4}
\opt{D}{$-4$}
\end{qoptions}
\qn{9} The value of `$r$' for which $x^{-17}$ is contained by the $r^{th}$ term in the expansion of $(x^4-x^{-3})^{15}$, is -
\begin{qoptions}
\opt{A}{10}
\opt{B}{11}
\opt{C}{12}
\opt{D}{13}
\end{qoptions}
\qn{10} $\displaystyle\sum_{r=0}^{n}r^2\left({}^nC_r\right)^2$ is equal to -
\begin{qoptions}
\opt{A}{$n^2.{}^{2n}C_n$}
\opt{B}{$n^2.{}^{2n}C_{n-1}$}
\opt{C}{$n^2.{}^{2n-2}C_n$}
\opt{D}{$n^2.{}^{2n-2}C_{n-1}$}
\end{qoptions}
\qn{11} If $\dfrac{{}^nC_r}{{}^nC_{(r+1)}}=\dfrac12$ and $\dfrac{{}^nC_{(r+1)}}{{}^nC_{(r+2)}}=\dfrac23$, then -
\begin{qoptions}
\opt{A}{$\dfrac nr=\dfrac72$}
\opt{B}{$\dfrac nr=7$}
\opt{C}{$\dfrac nr=\dfrac73$}
\opt{D}{$\dfrac nr=14$}
\end{qoptions}
\qn{12} Number of ways in which 10 examination papers can be arranged so that the best and the worst papers never come together, is -
\begin{qoptions}
\opt{A}{$2({}^9C_2)$}
\opt{B}{$2({}^9C_2)8!$}
\opt{C}{$2({}^9C_2)7!$}
\opt{D}{$8!$}
\end{qoptions}
\qn{13} LCM of 3 natural numbers $a,b,c$ is 150. If $N$ is the number of such ordered triplet, then -
\begin{qoptions}
\opt{A}{$N$ is multiple of 5}
\opt{B}{$N$ is an even number}
\opt{C}{$N$ is multiple of 7}
\opt{D}{$N$ is greater then 1000}
\end{qoptions}
\qn{14} If the equation $(p^2-3p+2)x^2+(p-1)^2x+(p^2-4p+3)=0$ is satisfied for more than two different values of $x$, then $p$ is equal to -
\begin{qoptions}
\opt{A}{1}
\opt{B}{2}
\opt{C}{1,2,3}
\opt{D}{2,3}
\end{qoptions}
\qn{15} If $\alpha$ and $\beta$ are the roots of equation $27x^2-4x+64=0$, then the value of $\left(\dfrac{\alpha^{2/3}}{\beta^{1/3}}+\dfrac{\beta^{2/3}}{\alpha^{1/3}}\right)$ is -
\begin{qoptions}
\opt{A}{$\dfrac13$}
\opt{B}{$\dfrac14$}
\opt{C}{$\dfrac19$}
\opt{D}{$\dfrac1{16}$}
\end{qoptions}
\qn{16} Sum of all the integers that do not belong to range of $y=\dfrac{11x^2+12x+6}{x^2+4x+2}$ is -
\begin{qoptions}
\opt{A}{$-7$}
\opt{B}{$-9$}
\opt{C}{$-4$}
\opt{D}{$-2$}
\end{qoptions}
\qn{17} Which of the following pair of function have same graph -
\begin{qoptions}
\opt{A}{$f(x)=\mathrm{sgn}\left(|x|+\dfrac1{|x|}\right)$; $g(x)=\left(\dfrac{x^2+1}{2|x|}\right)$}
\opt{B}{\raggedright $f(x)=\dfrac{(\sin x+\cos x)^2-1}{2\cos^2x}$; $g(x)=\dfrac{2\sin^2x}{1-(\sin x-\cos x)^2}$\par}
\opt{C}{$f(x)=e^{\ln(\mathrm{sgn}\,x^2)}$, $g(x)=\dfrac1{|\mathrm{sgn}(x)|}$}
\opt{D}{$f(x)=\sec^2x-\tan^2x$, $g(x)=1$\\
where [.] represents greatest integer function and sgn (.) represent singum function of $x$}
\end{qoptions}
\qn{18} Domain of function $f(x)=\log_2\sin(x)+\sqrt{4-x^2}$ is -
\begin{qoptions}
\opt{A}{(0, 2]}
\opt{B}{(1, 3]}
\opt{C}{(1, $\pi$]}
\opt{D}{[$-$2, 2]}
\end{qoptions}
\qn{19} Consider $f(x)=\dfrac1{(\sin x+2\cos x)(4\sin x-2\cos x)}$, then -
\begin{qoptions}
\opt{A}{range of $f(x)$ is $\left[-\dfrac1{\sqrt5},\ \dfrac1{\sqrt5}\right]$}
\opt{B}{domain of $f(x)$ is $R-(n\pi+\alpha)\ \forall\ \alpha=\tan^{-1}\left(\dfrac43\right)$, $n\in I$}
\opt{C}{domain of $f(x)$ is $R-\left(\dfrac{n\pi}2+\dfrac\alpha2\right)\ \forall\ \alpha=\tan^{-1}\dfrac43$, $n\in I$}
\opt{D}{range of $f(x)$ is $\left(-\infty,\ -\dfrac1{\sqrt5}\right]\cup\left[\dfrac1{\sqrt5},\ \infty\right)$}
\end{qoptions}
\qn{20} Which of the following best suits the graph of $f(x)=\sqrt{\mathrm{Sgn}|x|}$
\begin{qoptions}
\opt{A}{\qfigi{p2_math_s1_q20_A}}
\opt{B}{\qfigi{p2_math_s1_q20_B}}
\opt{C}{\qfigi{p2_math_s1_q20_C}}
\opt{D}{\qfigi{p2_math_s1_q20_D}}
\end{qoptions}
\sect{SECTION-II}
\qn{1} If $2C_0C_2+3C_1C_3+4C_2C_4+\dots+nC_{n-2}C_n=S(n)$ (where $C_r={}^nC_r$), then the value of $\dfrac{S(8)}{(220)(26)}$ is
\qn{2} If $1.\left({}^{20}C_0\right)^2+2.\left({}^{20}C_2\right)^2+3.\left({}^{20}C_4\right)^2+\dots+11.\left({}^{20}C_{20}\right)^2$ is equal to $3\left[{}^{8a}C_{4a}+{}^{4a}C_{2a}\right]$, then $a$ is equal to
\qn{3} Let $S=\displaystyle\sum_{n=1}^{\infty}\left(\dfrac{n^2}{\sum\limits_{r=0}^{n}r\left({}^nC_r\right)}\right)$, then the value of $S$ is
\qn{4} The postal stamps of three different valuations viz Rs 2, Rs 3 and Rs 4 are available. If the number of ways in which these stamps can be posted on envelop one after the other so that total valuation of stamps on envelop is Rs 10 is $N$, then $\left[\dfrac N2\right]$ is equal to (where [.] denotes greatest integer function)
\qn{5} If number of seven digit numbers in which first four digits are in ascending order, last four digits are in descending order and middle digit is a perfect square is $\lambda$, then digit at the tens place of $\lambda$ is
\qn{6} If the number of positive two digit integers which are divisible by both of its digit be $N$ then the value of $|N-10|$ is
\qn{7} Let there is a set of 12 natural numbers $A=\{481,482,\dots,492\}$. The number of subset $B$ of $A$ which contain at least 3 elements such that difference of any two elements is at least 2 is -
\qn{8} Let $x$ is the number of 6 digit even numbers such that sum of their digits is 50, then $\dfrac{x}{10}$ is equal to
\qn{9} If $a,b,c$ are the roots of the equation $x^3-x-1=0$, then the value of $\dfrac{(b+c)^2}{a^3-2}+\dfrac{(c+a)^2}{b^3-2}+\dfrac{(a+b)^2}{c^3-2}$ is
\qn{10} The smallest integral value of $a$ for which the inequality $1+\log_5(1+x^2)\le\log_5(ax^2+4x+a)$ is true for all $x\in R$ is
\end{multicols}
\end{document}
